\documentclass[../../../include/open-logic-section]{subfiles}

\begin{document}

\olfileid{sth}{ordinals}{replacement}
\olsection{પ્રતિસ્થાપન}

\olref[sth][ordinals][vn]{sec}માં આપણે ક્રમવાચક સંખ્યાઓ રજૂ કરવા માટે
એવો વિચાર કર્યો હતો કે તેમને ક્રમપ્રકારો, એટલે કે સુક્રમોના પ્રમાણભૂત
અવેજો, તરીકે લઈ શકાય. આ વિચાર સફળ થવા માટે આપણે સાબિત કરવું પડે કે
\emph{દરેક સુક્રમ કોઈ ક્રમવાચક સંખ્યા સાથે ક્રમ-એકરૂપ છે}.
પછી $\tuple{A, <} \isomorphic \alpha$ થાય એવી ક્રમવાચક
સંખ્યા $\alpha$ તરીકે $\ordtype{A, <}$ વ્યાખ્યાયિત કરી શકીએ.

દુર્ભાગ્યે, અત્યાર સુધી રજૂ કરેલાં સ્વયંસિદ્ધિઓથી જ આપણે આ ઇચ્છિત
પરિણામ \emph{સાબિત કરી શકતા નથી}. તેનું કારણ
\olref[replacement][strength]{sec}માં જોઈશું; હાલ એટલું જાણવું
પૂરતું છે કે આ પરિણામ મળતું નથી. તેથી એક નવો વિચાર જોઈએ:

\begin{axiom}[પ્રતિસ્થાપનનું પ્રરૂપ]
કોઈ પણ સૂત્ર $\phi(x, y)$ માટે નીચેનું એક સ્વયંસિદ્ધિ છે:
\begin{quote}
	કોઈ પણ $A$ માટે, જો $(\forall x \in A)\lexists![y][\phi(x,y)]$ હોય, તો
	$\Setabs{y}{(\exists x \in A)\phi(x,y)}$ અસ્તિત્વમાં છે.
\end{quote}
\end{axiom}
\noindent
પૃથક્કરણની જેમ આ પણ પ્રરૂપ છે: જુદાં જુદાં અસંખ્ય $\phi$ સૂત્રોમાંથી
દરેક માટે તે એક સ્વયંસિદ્ધિ આપે છે. તેને સમકક્ષ રીતે, અને સામાન્ય રીતે,
આ પ્રમાણે પણ લખાય છે:

\begin{defish}
``$B$'' ન આવતું હોય એવા કોઈ પણ સૂત્ર $\phi(x,y)$ માટે નીચેનું એક સ્વયંસિદ્ધિ છે:
\[
\forall A[(\forall x \in A)\lexists![y][\phi(x,y)] \lif \exists B\forall y (y \in B \liff (\exists x \in A)\phi(x,y))]
\]
\end{defish}

જોકે, પહેલી વાર જોતાં આ સૂત્ર ગૂંચવણભર્યું લાગે છે. પ્રતિસ્થાપનનું
આગળનું તાત્કાલિક પરિણામ આપણે સમજવા માગીએ છીએ તે વિચારને કદાચ વધુ
\emph{સ્પષ્ટ રીતે} વ્યક્ત કરે છે:\footnote{પદ એવી અભિવ્યક્તિ છે
જે તેના મુક્ત ચલોને મૂલ્યો આપીએ ત્યારે બરાબર એક જ વસ્તુ નિર્દેશે છે.
દાખલા તરીકે, ``$\emptyset$'' ખાલી ગણને નિર્દેશતું પદ છે;
``$\{x\}$'' એ $x$ના એકઘટકી ગણને નિર્દેશતું પદ છે, ભલે $x$ ગમે તે હોય;
અને ``$x \cup y$'' એ $x$ અને $y$ના યોગગણને નિર્દેશતું પદ છે,
ભલે તેઓ ગમે તે હોય.}

\begin{cor}
કોઈ પણ પદ $\tau(x)$ અને કોઈ પણ ગણ $A$ માટે નીચેનો ગણ અસ્તિત્વમાં છે:
\[
	\Setabs{\tau(x)}{x \in A} = \Setabs{y}{(\exists x \in A)y = \tau(x)}.
\]
\end{cor}

\begin{proof}
$\tau$ \emph{પદ} હોવાથી $\forall x \lexists![y][\tau(x) = y]$.
તેથી ખાસ કરીને $(\forall x \in A)\lexists![y][\tau(x) = y]$.
આથી પ્રતિસ્થાપન મુજબ
$\Setabs{y}{(\exists x \in A)\tau(x) = y}$ અસ્તિત્વમાં છે.
\end{proof}
\noindent
આ પરથી ``પ્રતિસ્થાપન'' નામ યોગ્ય લાગે છે: આપેલા ગણ $A$ના દરેક સભ્યને
$\tau$ હેઠળ તેના પ્રતિબિંબથી બદલીને નવો ગણ
$\Setabs{\tau(x)}{x \in A}$ બનાવી શકાય છે. ખરેખર, કોઈ વિધેય
હેઠળ $A$ના પ્રતિબિંબ માટેના સંકેત મુજબ, $\Setabs{\tau(x)}{x \in A}$ને
$\funimage{\tau}{A}$ પણ લખી શકીએ.

પરંતુ અહીં મહત્ત્વનું એ છે કે $\tau$ એક \emph{પદ} છે. તે
\olref[sfr][fun][rel]{sec}ના અર્થમાં, એટલે કે ક્રમિત જોડીઓના
કોઈ ચોક્કસ ગણ તરીકે, \emph{વિધેય} હોવું જરૂરી નથી, અને તેનું નામ
હોવું પણ જરૂરી નથી. જો $f$ એ અર્થમાં વિધેય હોય, તો
$\funimage{f}{A} = \Setabs{f(x)}{x \in
A}$ એ માત્ર $\ran{f}$નો એક ઉપગણ છે. આ ઉપગણનું અસ્તિત્વ
તો $\Zminus$નાં સ્વયંસિદ્ધિઓથી જ મળે છે.\footnote{માત્ર
$\Setabs{y \in \bigcup \bigcup f}{(\exists x \in A)y =
f(x)}$ લો.} તેના મુકાબલે, પ્રતિસ્થાપન આપણા સ્વયંસિદ્ધિઓમાં
\emph{શક્તિશાળી} ઉમેરો છે, જે
\olref[sth][replacement][]{chap}માં જોઈશું.


\end{document}
